\subsection{Physics-Informed Neural Networks}
\label{sec:physics-informed}

Physics-informed neural networks (PINNs) embed physical laws as constraints in the loss function, ensuring that predictions respect known physical behavior \cite{}. This approach has shown promise in structural health monitoring, where domain-specific knowledge can guide the learning process.

The degradation-trend-constrained variational autoencoder (DTC-VAE) represents a notable advancement in physics-constrained deep learning for guided wave SHM \cite{dtcvae2026}. By incorporating a monotonicity constraint into the VAE latent space, DTC-VAE ensures that the learned representations encode physically meaningful degradation trends. Similarly, energy-oriented autoencoders have been developed for health state assessment, leveraging energy conservation principles to improve feature interpretability \cite{perry2026health}.

These physics constraints improve feature preservation by encoding domain-specific invariants directly into the network architecture or optimization objective. Rather than relying solely on data-driven learning, physics-informed models can maintain acoustically meaningful characteristics such as arrival times, amplitudes, and waveform morphology even under noisy conditions.

Despite these advances, physics constraints remain rarely applied to laser ultrasonic denoising. Existing deep learning approaches for laser US typically employ generic denoising objectives without explicit physical regularisation. This gap presents an opportunity: embedding physics-based loss terms that preserve acoustic features could yield denoised signals that better support defect detection while reducing reliance on extensive signal averaging.

% word count: ~190